\documentclass{article}
\usepackage[T1]{fontenc}
\usepackage{lmodern}
\usepackage{graphicx}
\usepackage{amsmath,amssymb}
\usepackage[a4paper,margin=2cm]{geometry}
\usepackage[absolute,overlay]{textpos}
\setlength{\TPHorizModule}{1mm}
\setlength{\TPVertModule}{1mm}
\usepackage[indonesian]{babel}
\usepackage{hyperref}
\setlength{\emergencystretch}{2em}
% unit: Halaman 11

\begin{document}

% === Halaman 11 (python/native) ===

11

1. Electronic Code Book (ECB)

• Setiap blok plainteks Pi dienkripsi secara individual dan independen 

dari blok lainnya menjadi blok cipherteks Ci . 

• Enkripsi: Ci = EK(Pi)

 Dekripsi: Pi = DK(Ci)

 yang dalam hal ini, Pi dan Ci masing-masing blok plainteks dan 

cipherteks ke-i.

\end{document}
